\documentclass[12pt]{article}

% a template that a friend gave, it's worked well enough for me
% i have added some packages and stuff that have proved useful

\usepackage{fancyhdr}
\usepackage{tipa}
\usepackage{fontspec}
\usepackage{amsfonts}
\usepackage{enumitem}
\usepackage[margin=1in]{geometry}
\usepackage{graphicx}
\usepackage{float}
\usepackage{amsmath}
\usepackage{braket}
\usepackage{amssymb}
\usepackage{booktabs}
\usepackage{hyperref}
\usepackage{mathtools}
\usepackage{xcolor}
\usepackage{float}
\usepackage{algpseudocodex}
\usepackage{titlesec}
\usepackage{bbm}

\pagestyle{fancy}
\fancyhf{} % sets both header and footer to nothing
\lhead{Kevin Sheng}
\setmainfont{Comic Neue}
\renewcommand{\headrulewidth}{1pt}
\setlength{\headheight}{0.75in}
\setlength{\oddsidemargin}{0in}
\setlength{\evensidemargin}{0in}
\setlength{\voffset}{-.5in}
\setlength{\headsep}{10pt}
\setlength{\textwidth}{6.5in}
\setlength{\headwidth}{6.5in}
\setlength{\textheight}{8in}
\renewcommand{\headrulewidth}{0.5pt}
\renewcommand{\footrulewidth}{0.3pt}
\setlength{\textwidth}{6.5in}
\usepackage{setspace}
\usepackage{multicol}
\usepackage{float}
\setlength{\columnsep}{1cm}
\setlength\parindent{24pt}
\usepackage [english]{babel}
\usepackage [autostyle, english = american]{csquotes}
\MakeOuterQuote{"}

\setlength{\parskip}{6pt}
\setlength{\parindent}{0pt}

\titlespacing\section{0pt}{12pt plus 4pt minus 2pt}{0pt plus 2pt minus 2pt}
\titlespacing\subsection{0pt}{12pt plus 4pt minus 2pt}{0pt plus 2pt minus 2pt}
\titlespacing\subsubsection{0pt}{12pt plus 4pt minus 2pt}{0pt plus 2pt minus 2pt}

\hypersetup{colorlinks=true, urlcolor=blue}

\newcommand{\correction}[1]{\textcolor{red}{#1}}

\usepackage{fancyvrb}

\rhead{Math 168}

\DeclareMathOperator{\Fr}{Fr}
\DeclareMathOperator{\Tr}{Tr}
\newcommand{\lra}{\xLeftrightarrow}
\newcommand{\ra}{\xRightarrow}

\setlist[enumerate,1]{label=(\alph*)}

\begin{document}

\section{Bipartite Distributions}

For convenience, I'll represent the degree distributions as
$a_0, a_1, \cdots$ and $b_0, b_1, \cdots$ instead, where $a_i$ is the chance
of a node in $A$ having degree $i$ and similarly for $b_i$.

\subsection{Part A}

\boxed{g_0(z)=\sum_{i=0}^{\infty} a_iz^i} and
\boxed{h_0(z)=\sum_{i=0}^{\infty} b_iz^i}.

Getting $g_1$ and $h_1$ is a bit trickier, though.
Since the graph is bipartite, to get to any $a \in A$, we have to follow it from some $b \in B$.

If we fix a stub off of $b$, then the chance it links to a specific stub
with degree $k$ is $\frac{k}{m}$.
There's around $a_k|A|$ total nodes with degree $m$, so
\[P(\text{our stub ends at a node with degree $k$})=\frac{k}{m} \cdot a_k|A| = \frac{ka_k}{\braket{a}}\]
where $\braket{a}$ is the average degree of all nodes in $A$.

Thus, \boxed{g_1(z)=\sum_{i=0}^{\infty} \frac{(i+1)a_{i+1}}{\braket{a}}z^i},
and by symmetry \boxed{h_1(z)=\sum_{i=0}^{\infty} \frac{(i+1)b_{i+1}}{\braket{b}}z^i}.

\subsection{Part B}

In the 1-mode projection onto $A$, $a_1, a_2 \in A$ are connected iff $\exists b \in B$
that's connected to both of them.
This means that we can go to all neighbors in $b$ and tally all the other $a$s
(minus ourselves) that they're connected to get $a$'s degree in the projection.
Notice that what I just described is literally just the number of second neighbors
that $a$ has.

Thus, $G_k$, the chance $a$ has degree $k$ in the projection, is
\[G_k = \sum_{i=0}^{\infty} a_i P[\text{$a$ has $k$ second neighbors} \mid \deg(a)=i]\]
I'll abbreviate the probability as $P[k \mid i]$ from here on out.

Now let's do some work on the PGF:
\begin{align*}
      G(z)
       & = \sum_{i=0}^{\infty} G_i z^i                                \\
       & = \sum_{i=0}^{\infty} \sum_{k=0}^{\infty} a_k P[i \mid k]    \\
       & = \sum_{k=0}^{\infty} a_k \sum_{i=0}^{\infty} P[i \mid k]z^i
\end{align*}
By the identities and arguments proven in class,
\[\sum_{i=0}^{\infty} P[i \mid k]z^i = h_1(z)^k\]
and we can continue:
\begin{align*}
      G(z)
       & = \sum_{k=0}^{\infty} a_k \sum_{i=0}^{\infty} P[i \mid k]z^i \\
       & = \sum_{k=0}^{\infty} a_k h_1(z)^k                           \\
       & = g_0(h_1(z))\quad\square
\end{align*}

\subsection{Part C}

The expected value given the PGF $z$ can be found by evaluating $G'(1)$, which is
\begin{align*}
      G'(1)
       & = g_0'(h_1(1)) \cdot h_1'(1)                                                \\
       & = g_0'(1)h_1'(1)                                                            \\
       & = \boxed{\frac{\braket{a}\left(\Braket{b^2}-\braket{b}\right)}{\braket{b}}}
\end{align*}

\pagebreak

\section{Stochastic Probabilities}

\begin{enumerate}
      \item There's $\binom{n_1-1}{k}$ ways for a node in the first group
            to connect to $k$ others of its own kind, so the chance is
            \boxed{\binom{n_1-1}{k}p_1^k(1-p_1)^{n_1-1-k}}.

      \item By similar reasoning, we have \boxed{\binom{n_2}{k}q^k(1-q)^{n_2-k}}.

      \item We have the PMFs of the number of connections to the first and second block.
            If we let these be represented by random variables $X$ and $Y$ respectively, we have
            \begin{align*}
                  P(\deg v = k)
                   & = \sum_{i=0}^{k} P(X=i)P(Y=k-i)                                                                        \\
                   & = \sum_{i=0}^{k} \binom{n_1-1}{i}p_1^i(1-p_1)^{n_1-1-i} \cdot \binom{n_2}{k-i}q^{k-i}(1-q)^{n_2-(k-i)}
            \end{align*}
            Yeah...I don't see any elegant way to reduce this.

      \item There's $n_2$ chances to connect to a neighbor in the second group,
            and chance works with probability $q$, so the expected
            number of neighbors in the second group is \boxed{qn_2}.
\end{enumerate}

\section{Watts-Strogatz Graphs}

See the attached Python notebook.

\end{document}
